\section{Goodness-of-Fit testing : Definitions and criteria}


\section{Bayesian approach for goodness-of-fit testing}
One can note that since the set of admissible test is too wide, there are no general recommendations as to how to select optimal tests when the alternative is composite and of dimension more than one.
A way out of this situation is to characterize the test, not in the usual way by using error functions $\alpha(\psi,P),\beta(\psi,Q)$ for some $P\in \mp_0, Q\in\mp_1$, but by some monotone functional $F_1(\psi),F_2(\psi)$
\subsection{General setting}
Fix a Borel $\sigma$-algebra on the space $\mp$ of probability measures (with respect to the total variation distance) such that $\mp_0,\mp_1$ are Borel subsets. Moreover, fix two priors $\pi_0,\pi_1$ with supports in $\mp_0,\mp_1$, respectively. \\
Now, define two functionals $F_1,F_2$ as the average errors with respect to each prior : 
\begin{equation}\label{bayesdef}
    \begin{aligned}
         F_1(\psi) &\coloneqq \alpha(\psi, \pi_0) = \E_{\pi_0} \alpha(\psi,P) = \int \alpha(\psi,P)d\pi_0(P)\\
         F_2(\psi) &\coloneqq \beta(\psi,\pi_1) = \E_{\pi_1} \beta(\psi,P) = \int \beta(\psi,P)d\pi_1(P)
    \end{aligned}
\end{equation}
Analogously to what has been done with the standard approach, given $\alpha \in \left(0,1\right)$ we define the set : 
\begin{equation}\label{setprior}
\Psi_{\alpha,\pi_0} \coloneqq \bigl\{\psi \in \Psi : \alpha(\psi,\pi_0) \leq \alpha \bigr\}
\end{equation}
and the following functions : 
\begin{align}
    \gamma_t(\psi;\pi_0,\pi_1) &\coloneqq t\alpha(\psi,\pi_0) + \beta(\psi,\pi_1) \notag\\
    \gamma(t;\pi_0,\pi_1) &\coloneqq \inf_{\psi \in \Psi} \gamma_t(\psi;\pi_0,\pi_1) \label{gammabayes}\\
    \beta(\alpha;\pi_0,\pi_1) &\coloneqq \inf_{\psi \in \Psi_{\alpha,\pi_0}} \beta(\psi,\pi_1) \label{betabayes}
\end{align}
The tests $\psi^{(t)}$ or $\psi_\alpha$ are called Bayesian if they provide an infimum \cref{gammabayes} or in \cref{betabayes}. 
\begin{definition}[Mixture of a prior] 
    Given a prior $\pi$, one can define its mixture as the measure $P_\pi$ on $(\mx,\mathcal{A})$ given by : 
    \begin{equation*}
        \forall A \in \mathcal{A} : \P_\pi (A) = \E_\pi P(A) = \int_\mp P(A) d\pi(P)
    \end{equation*}
\end{definition}

Given all these definitions, the Bayesian hypothesis testing is reduced to testing of the simple hypotheses with null hypothesis $P = P_{\pi_0}$ and alternative $P = P_{\pi_1}$
\begin{remark}
    Note that when we have a parametrization which corresponds to an injective map $\Theta \rightarrow \mp$, we set a Borel $\sigma$-algebra $\mathcal{T}$ on $\Theta$ (with respect to some topology) that makes the map measurable, one can assume that the priors $\pi_0,\pi_1$ are defined on ($\Theta,\mathcal{T}$). Hence, in our specific context, one can assume that $\pi_0,\pi_1$ are measures on ($\R^n,\mathcal{B}^n$) in the finite-dimensional case, or on ($\ell^2,\mathcal{B}$) in the infinite-dimensional case (Here $\mathcal{B}$ denotes a Borel $\sigma$-algebra on $\ell^2$).
\end{remark}
\textcolor{red}{Add details + define the following measure}
\begin{equation}\label{measureanalytic}
    r(dv) = C^{-1}\pi(dv)e^{-\frac{\left|v\right|^2}{2}} , C = \int e^{-\frac{\left|v\right|^2}{2}} \pi(dv) = L_\pi(0)
\end{equation}
\begin{prop}\label{prop2.1}
    Let $\pi_0,\pi_1$ be two priors. Then the following equivalence holds : 
    \begin{equation*}
        P_{\pi_0} = P_{\pi_1} \Leftrightarrow \pi_0 = \pi_1
    \end{equation*}
\end{prop}
\begin{proof}
    Clearly, if $\pi_0 = \pi_1$, then $P_{\pi_0} = P_{\pi_1}$. Now, assume that $P_{\pi_0} = P_{\pi_1}$ then $L_{\pi_0} = L_{\pi_1}$ almost surely. 
    First, consider the finite-dimensional case $x \in \R^n$. Let $\mu_1,\mu_2$ be the probability measures defined in \cref{measureanalytic} for the priors $\pi_0,\pi_1$. By the \textbf{previous} proposition, the functions $L_{\pi_0},L_{\pi_1}$ are analytical, this implies that the equality $L_{\pi_0}=L_{\pi_1}$ holds true in the whole space $\mathbb{C}^n$. In particular, we get the equality for $z = ix, x\in \R^n$ which corresponds to the equality of characteristic functions of the probability measures $\mu_0,\mu_1$. Therefore $\mu_0 = \mu_1$ which yields $\pi_0 = \pi_1$. 
    In the infinite-dimensional case, the proof above yields equality for all finite dimensional marginals, which implies $\pi_0 = \pi_1$.
\end{proof}


\section{Minimax approach in Hypothesis testing}
\subsection{General setting}
Instead of considering the average error, we consider maximum error, in this section, define : 
\begin{align}
    F_1(\psi) = \alpha(\psi,\mp_0) &= \sup_{P\in\mp_0}\alpha(\psi,P) \\
    F_2(\psi) = \beta(\psi,\mp_1) &= \sup_{P\in\mp_1}\beta(\psi,P)
\end{align}
Moreover, similarly to \cref{setprior}, given $\alpha \in \left(0,1\right)$ we define : 
\begin{equation}
    \Psi_\alpha = \Psi_{\alpha,\mp_0} = \bigl\{\psi \in \Psi : \alpha(\psi,\mp_0) \leq \alpha \bigr\}
\end{equation}
One can also define the following quantities : 
\begin{align}
    \gamma_t(\psi,\mp_0,\mp_1) &= t\alpha(\psi,\mp_0) + \beta(\psi,\mp_1) \notag \\
    \gamma(t,\mp_0,\mp_1) &= \inf_{\psi \in \Psi} \gamma_t(\psi) \label{minimaxt} \\
    \beta(\alpha,\mp_0,\mp_1)&= \inf_{\psi\in \Psi_{\alpha,\mp_0}} \beta(\psi,\mp_0,\mp_1)  \label{minimaxalpha}
\end{align}
The tests $\psi^{(t)}$ or $\psi_\alpha$ are called \textit{minimax} if they provide infimums for \cref{minimaxt} or \cref{minimaxalpha}. \\ 
\newpage
One can observe the following convexity lemma : 
\begin{lemma}\label{conv}
    The following properties hold true : 
    \begin{enumerate}
        \item The functions $\alpha(\psi,\mp_0),\beta(\psi,\mp_1),\gamma_t(\psi,\mp_0,\mp_1)$ are convex in $\psi\in\Psi$
        \item The function $\alpha \mapsto \beta(\alpha)$ is convex nonincreasing continuous for $\alpha \in [0,1]$, and the function $t\mapsto \gamma(t)$ is concave nondecreasing continuous in $t \in [0,\infty )$. These two functions are connected by the relations : 
        \begin{equation}\label{gammabetarelation}
            \gamma(t) = \inf_{\alpha \in [0,1]} \big(t\alpha + \beta(\alpha)\big) \text{ and } \beta(\alpha) = \sup_{t \in [0,\infty)} \big(\gamma(t) - t\alpha\big)
        \end{equation}
    \end{enumerate}
\end{lemma}
\textcolor{red}{Maybe add a proof}
\subsection{Connection with Bayesian approach}
Now that the settings of both Minimax and Bayesian approaches were explored, we will see that there exists a close relation between these two. All the following statements go back to \citet{Wald1950-WALSDF-3}
\begin{definition}[$L_1$ distance between two families of measures]
    Let $A,B$ be two families of measures, we define the $L_1$ distance between $A$ and $B$ as follows : \begin{equation*}
        d(A,B) \coloneqq \inf_{P\in A, Q \in B} TV(P,Q)
    \end{equation*}
    Moreover, we define $[A]$ to be the convex hull of $A$.
\end{definition}

\begin{theorem}\label{thm2.1}
    \textcolor{red}{Add some text}
    \begin{enumerate}
        \item The following inequalities hold true : 
        \begin{align}
            \gamma(t) &\geq \sup \bigl\{\gamma(t,P_0,P_1) : P_0 \in [\mp_0], P_1 \in [\mp_1] \bigr\} \label{ineqgamma}\\
            \beta(\alpha) &\geq \sup \bigl\{\beta(\alpha, P_0, P_1) : P_0 \in [\mp_0], P_1 \in [\mp_1]\bigr\} \label{ineqbeta}
        \end{align}
        In particular, $\gamma = \gamma(1) \geq 1-\frac{1}{2}d([\mp_0],[\mp_1])$.
        \item If $\mp_0,\mp_1$ are dominated families of measures on $(\mx,\mathcal{A})$, and the $\sigma-$algebra $\mathcal{A}$ is countably generated, then for any $t>0, \alpha \in \left(0,1\right)$ the infimas in $\ref{ineqgamma}$ and $\ref{ineqbeta}$ are attained, so the inequalities become equalities.
        \item If $\mp_0,\mp_1$ are compact with respect to the total varation distance. Then, there exists pairs of priors $\pi_{0,t},\pi_{1,t}$ and $\pi_{0,\alpha},\pi_{1,\alpha}$ such that the minimax tests $\psi^{(t)}, \psi_\alpha$ are optimal tests for testing the simple hypothesis $P_{\pi_{0,t}}$ or $P_{\pi_{0,\alpha}}$ against the simple alternative $P_{\pi_{1,t}}$, or $P_{\pi_{1,\alpha}}$. The priors $\pi_{0,t},\pi_{1,t}$ and $\pi_{0,\alpha},\pi_{1,\alpha}$ are called least favourable.
    \end{enumerate}
\end{theorem}
\begin{proof}
    \textcolor{red}{ADD PROOF}
\end{proof}
We will use the previous inequalities \ref{ineqgamma},\ref{ineqbeta} in order to construct some lower bounds for the quantities $\gamma(t), \beta(\alpha),\gamma$. For example, let $\pi_{0,n},\pi_{1,n}$ be sequences of priors such that : 
\begin{equation}\label{condpriors}
    \pi_{0,n}(\mp_0) = \pi_{1,n}(\mp_1) = 1
\end{equation}
Then, one can rewrite the inequalities as : 
\begin{align}
    \gamma(t) \geq \limsup_{n \rightarrow\infty} \gamma(t;P_{\pi_{0,n}},P_{\pi_{0,n}})& \label{limsupgamma} \\
    \beta(\alpha) \geq \limsup_{n\rightarrow \infty} \beta(\alpha; P_{\pi_{0,n}},P_{\pi_{1,n}})& \label{limsupbeta} \\ 
    \gamma \geq 1-\frac{1}{2}\liminf_{n\rightarrow \infty} d(P_{\pi_{0,n}},P_{\pi_{1,n}})& \label{limsupgamma1}
\end{align}
One can see that this can be also true if $ \pi_{0,n}(\mp_0) \rightarrow 1$ and $ \pi_{1,n}(\mp_1) \rightarrow 1$ 
\begin{prop}
    Let $\pi_{0,n},\pi_{1,n}$ be sequences of priors such that : 
    \begin{equation*}
        \lim_{n\rightarrow \infty} \pi_{0,n} (\mp_0) = 1, \lim_{n\rightarrow \infty} \pi_{1,n} (\mp_1) = 1 
    \end{equation*}
    Then, the inequalities \ref{limsupgamma}, \ref{limsupbeta} and \ref{limsupgamma1} hold true. 
\end{prop}
\begin{proof}
    \textcolor{red}{Add proof prop 2.9 IngSus}
\end{proof}
\begin{theorem}\citet[Theorem 2.3]{Ing2012}\\
    Let $\pi_0,\pi_1$ be priors, and let $S_0 \subseteq \mp_0, S_1 \subseteq \mp_1$ be measurable sets such that $\pi_0(S_0) = \pi_1(S_1) = 1$. Moreover, assume that the Bayesian tests $\psi^{t}$ or $\psi_{\alpha}$ satisfy for $P_0 \in \mp_0,P_1 \in \mp_1$ : 
    \begin{align}
        \alpha(\psi_\alpha,\mp_0) = \alpha(\psi_\alpha,P_0) \ &, \ \beta(\psi_\alpha,\mp_1) = \beta(\psi_\alpha,\mp_1) \text{ or } \notag \\
        \gamma_t(\psi^{(t)};\mp_0,\mp_1) &= \psi(\psi^{(t)};P_0,P_1) \notag
    \end{align}
    Then, $(\pi_0,\pi_1)$ is a pair of least favourable priors under the criteria $\gamma(t)$ or $\beta(\alpha)$, and the tests $\psi^{(t)}$ or $\psi_\alpha$ are minimax. 
\end{theorem}
\begin{proof}
    \textcolor{red}{Add proof .}
\end{proof}